\chapter{Funktionsweise des Programms im Überblick}
\label{ch:funktionsweise}

Dieses Kapitel zeigt kurz, wie das Programm als Ganzes funktioniert. Die folgenden Kapitel gehen dann ins Detail.

\section{Grundidee}

Das Programm beantwortet eine Frage: \emph{Wie viel Energie kostet eine echte Fahrt ein E-Bike, und wie unterscheiden sich dabei zwei Akkutypen?} Dazu geht es die GPS-Aufzeichnung Punkt für Punkt durch und berechnet daraus Kraft, Leistung, Drehmoment und Motorstrom. Dieser Strom wirkt dann auf ein Akkumodell, das Ladezustand und Spannung entsprechend fortschreibt.

Die Bewegung (Geschwindigkeit, Beschleunigung, Steigung) liegt bereits als GPS-Messung vor. Aus dieser bekannten Bewegung berechnet das Programm über das Freikörperdiagramm die dafür nötige Kraft und daraus den Strom.

\section{Ablauf in neun Schritten}

\begin{enumerate}
  \item \textbf{CSV einlesen} (\texttt{gps\_track.py}): Die Datei \texttt{data/final\_project\_input\_data.csv} wird eingelesen. Jede Zeile enthält Zeit, Position, Höhe und Temperatur.
  \item \textbf{Kinematik berechnen}: Aus je zwei Punkten entstehen Zeitdifferenz, Distanz (Haversine-Formel), Rohgeschwindigkeit, Steigung, Kurswinkel und Kompassrichtung.
  \item \textbf{Geschwindigkeit optional glätten} (\texttt{speed\_smoothing.py}): Kurze Zeitabstände lassen die Rohgeschwindigkeit stark schwanken. Ist die Glättung aktiviert, dämpft ein Medianfilter mit anschließender Exponentialglättung diese Spitzen. Die Beschleunigung wird danach aus der aktiven Geschwindigkeitskurve berechnet.
  \item \textbf{Luftdichte schätzen} (\texttt{luftdichte.py}): Aus Höhe und Temperatur ergibt sich die aktuelle Luftdichte, statt pauschal mit einem Standardwert zu rechnen.
  \item \textbf{Kraft, Leistung und Drehmoment berechnen} (\texttt{ebike.py}): Für jeden Punkt addiert \texttt{EBike.punkt\_auswerten()} Luftwiderstand, Hangabtrieb, Beschleunigungskraft und Rollwiderstand. Daraus folgen Leistung und Drehmoment.
  \item \textbf{Motorstrom bestimmen} (\texttt{motor.py}): Das Drehmoment wird über die Motorkonstante linear in einen Strom umgerechnet.
  \item \textbf{Akku fortschreiben} (\texttt{battery\_pack.py}, \texttt{lipo\_battery.py}, \texttt{nmc\_battery.py}, \texttt{akkutemperatur.py}): Der Strom wirkt über die Zeitspanne auf den Ladezustand. Die Temperatur korrigiert dabei Innenwiderstand und Kapazität, die Spannung ergibt sich aus der Leerlaufspannungskennlinie. Würde der Akku über 100\,\% geladen (z.\,B. bei starkem Bergabfahren), begrenzt \texttt{bremswiderstand.py} den Strom und wandelt den Rest in Wärme um.
  \item \textbf{Simulation wiederholen} (\texttt{ebike\_simulator.py}, \texttt{main.py}): Dieselbe Strecke läuft einmal mit LiPo- und einmal mit NMC-Akku, da beide von derselben Basisklasse erben.
  \item \textbf{Ausgeben} (\texttt{plot\_utils.py}, \texttt{gps\_track.py}): Am Ende stehen Kennzahlen, HTML-Karten und PNG-Diagramme zur Verfügung.
\end{enumerate}

\section{Zusammenspiel der Bausteine}

\texttt{main.py} ist der Einstiegspunkt: Es lädt Konfiguration und GPS-Daten und erzeugt die Objekte \texttt{GPSTrack}, \texttt{EBike}, \texttt{Motor} und zwei Akkus. \texttt{GPSTrack} liefert nur Bewegungsdaten, \texttt{EBike} nur die Fahrzeugphysik, \texttt{Motor} nur die Umrechnung von Drehmoment in Strom, die Akkuklassen nur ihr elektrisches Verhalten. \texttt{EBikeSimulator} verbindet diese vier Objekte, enthält aber selbst keine Physik -- er steuert nur den Ablauf und sammelt die Ergebnisse. Diese klare Trennung ist die Grundlage der Architektur aus Kapitel~\ref{ch:softwarearchitektur}.

\section{Die zwei Konfigurationsdateien}

Zwei YAML-Dateien steuern das Verhalten, ohne dass Code geändert werden muss: \texttt{data/bike\_config.yaml} legt die Fahrzeugparameter fest (Tabelle~\ref{tab:ebike-parameter}), \texttt{data/speed\_smoothing\_config.yaml} die Glättung (Tabelle~\ref{tab:smoothing-config}). Beide werden beim Start einmal geladen.

\section{Ergebnis in einem Satz}

Aus einer GPS-Aufzeichnung entsteht so ein nachvollziehbarer Weg von Position und Höhe über Kraft und Leistung bis zum Ladezustand zweier Akkus -- die Zahlen dazu stehen in Kapitel~\ref{ch:ergebnisse}.
